\chapter{Experimental Setup}
\label{chap:experimental-setup}

\section{Overview}
\label{sec:exp-overview}

This chapter defines the experimental methodology used to evaluate the
Monte Carlo benchmarking framework.  The aim is to ensure that all
measurements reported in Chapter~\ref{chap:results} are produced from
controlled, reproducible, and scientifically interpretable experiments.

The evaluation combines two forms of evidence. First, a detailed local CPU
campaign is performed on a GitHub Codespace with 4 CPU cores and 16 GB RAM.
This campaign supports the detailed validation, scalability, AD-overhead,
thread-scaling, memory-pressure, and workload-sensitivity analysis. Second, a
final cloud campaign evaluates cost-performance and profiler quality on Google
Cloud instances, with the headline \texttt{sha\_cloud\_profiler\_v4} run
executed on \texttt{n2-standard-4} in \texttt{europe-west1-b}. GPU execution,
distributed execution, mixed precision, and Mojo-based vendor-claim validation
remain future extensions rather than completed experiments.

The experimental design follows three principles.

\begin{enumerate}
    \item \textbf{Like-for-like workloads:} every engine is evaluated on the
    same workload configuration, random seed policy, path count, model
    parameters, and timing procedure.

    \item \textbf{Separation of concerns:} performance differences are
    isolated by varying one main factor at a time, such as engine, AD mode,
    path count, time-step count, or requested thread count.

    \item \textbf{Auditable measurement:} each benchmark run stores timing
    statistics, numerical outputs, memory measurements, software versions,
    hardware metadata, and the full workload configuration in the shared
    SQLite database.
\end{enumerate}

Three workloads are used across the complete evaluation. The European GBM
workload is used as a simple analytical validation case and low-cost baseline.
The parametric local-volatility workload is used as the main local performance
workload because it contains time stepping, a nonlinear volatility surface, and
a larger computational graph for automatic differentiation. The Asian option is
used in the final cloud-profiler experiments as a path-dependent workload that
exposes a different backend ranking from the two European-style workloads.

\section{Connection to the Research Questions}
\label{sec:exp-rq-mapping}

Table~\ref{tab:exp-rq-mapping} maps the project research questions to the
experiments that generate the evidence used in the results chapter. Cloud
hardware is evaluated through the final cost-performance and profiler
experiments. GPU acceleration, distributed execution, and vendor-specific Mojo
claims are not directly evaluated; they are discussed as scope limitations and
future work.

\begin{table}[htbp]
\centering
\small
\caption{Mapping from research questions to local experiment evidence.}
\label{tab:exp-rq-mapping}
\begin{tabularx}{\textwidth}{p{0.28\textwidth}Y p{0.25\textwidth}}
\toprule
\textbf{Research question} &
\textbf{Evidence collected} &
\textbf{Script(s)} \\
\midrule

How do software stacks affect MC+AD throughput and efficiency? &
Runtime, throughput, numerical output, and memory measurements are compared
across NumPy, JAX/XLA, C++/OpenMP, and Rust/Rayon under the same workload
definitions. &
\code{run_european_baseline.py};
\code{run_european_scalability.py};
\code{run_localvol_suite.py} \\

What are the trade-offs between programming languages and AD libraries? &
No-AD runtime and throughput are compared across Python/NumPy, JAX/XLA,
C++/OpenMP, and Rust/Rayon.  Forward- and reverse-mode AD are evaluated
through JAX, which provides the AD-capable reference implementation. &
\code{run_european_baseline.py};
\code{run_european_ad_analysis.py};
\code{run_localvol_suite.py} \\

How do forward- and reverse-mode AD affect MC workloads? &
AD overhead is measured by comparing JAX no-AD, forward-mode AD, and
reverse-mode AD runs at matched path counts, time-step counts, warmup
counts, and repetition counts. &
\code{run_european_ad_analysis.py};
\code{run_localvol_suite.py} \\

How do path count and time-step count affect scalability? &
The experiments sweep the number of Monte Carlo paths \(M\) and, for the
local-volatility workload, the number of time steps \(N\).  This separates
path-level scaling from discretisation-cost scaling. &
\code{run_european_scalability.py};
\code{run_localvol_suite.py} \\

How do parallelisation strategies affect scaling and bottlenecks? &
Strong and weak thread-scaling experiments vary the requested CPU thread
count over \(1,2,4\) threads.  Observed thread counts are stored where
available to distinguish requested from actual runtime behaviour. &
\code{run_thread_scalability.py};
\code{run_localvol_suite.py} \\

Which Google Cloud configurations provide good cost/performance? &
Runtime and estimated cost per run are compared for priced cloud instances,
using stored instance metadata and generated report tables. &
\code{run\_profiler\_vs\_grid.py};
\code{generate\_report\_assets.py} \\

Can profiling reduce full cloud benchmarking cost? &
The old profiler and SHA are compared with a full-grid baseline using full
runs saved, Pareto recovery, runtime regret, cost regret, and rank
correlation. &
\code{run\_profiler\_vs\_grid.py} \\

How far can the framework support broader proposal goals? &
The framework records enough configuration and environment metadata to
support cloud, GPU, distributed, mixed-precision, BLAS, or Mojo experiments,
although not all of these extensions are completed in the final evaluation. &
All experiment scripts and database records \\
\bottomrule
\end{tabularx}
\end{table}

\section{Execution Environment}
\label{sec:exp-environment}

The detailed local experiments are executed in a GitHub Codespace with 4 CPU
cores and 16 GB RAM.  The exact CPU model, operating system, Python version,
NumPy version, JAX version, and other software metadata are captured
automatically in the benchmark database for every run.

The use of a single local environment is deliberate: it reduces variation from
uncontrolled hardware changes and allows the local experiments to focus on the
framework, workload definitions, engines, AD modes, and CPU scaling behaviour.
Cloud conclusions are drawn only from the later cloud-profiler experiments and
are tied to their recorded instance types, regions, pricing assumptions, and
software environment.

\section{Workloads}
\label{sec:exp-workloads}

\subsection{European Option under Geometric Brownian Motion}
\label{sec:exp-workload-european}

The validation workload prices a European call option under geometric
Brownian motion (GBM).  Under the risk-neutral measure, the terminal asset
price is

\[
    S_T
    =
    S_0
    \exp\!\left[
        \left(r - \frac{1}{2}\sigma^2\right)T
        + \sigma\sqrt{T}Z
    \right],
    \qquad
    Z \sim \mathcal{N}(0,1).
\]

The Monte Carlo estimator for the discounted call price is

\[
    \hat{P}
    =
    e^{-rT}
    \frac{1}{M}
    \sum_{i=1}^{M}
    \max(S_T^{(i)} - K, 0),
\]

where \(M\) is the number of simulated paths.  Unless otherwise stated, the
fixed parameters are

\[
    S_0 = 100,\quad
    K = 100,\quad
    r = 0.05,\quad
    \sigma = 0.20,\quad
    T = 1.0,\quad
    \text{seed} = 42.
\]

This workload is used for analytical validation, simple baseline
performance, path-count scalability, and JAX AD correctness.  Its main
advantage is that the Black--Scholes analytical price and Greeks provide a
direct reference for the Monte Carlo price, Delta, Vega, and Rho.

\subsection{European Option under Parametric Local Volatility}
\label{sec:exp-workload-localvol}

The main performance workload prices the same European payoff, but replaces
constant volatility with a parametric local-volatility surface.  The asset
path is advanced using a log-Euler discretisation over \(N\) time steps:

\[
    S_{t+\Delta t}
    =
    S_t
    \exp\!\left[
        \left(r - \frac{1}{2}\sigma(S_t,t)^2\right)\Delta t
        + \sigma(S_t,t)\sqrt{\Delta t}Z_t
    \right],
    \qquad
    Z_t \sim \mathcal{N}(0,1).
\]

The local volatility surface is parameterised as

\[
    \sigma(S,t)
    =
    \sigma_{\min}
    +
    \operatorname{softplus}
    \left(
        a_0
        + a_1 \log(S/S_0)
        + a_2 \log(S/S_0)^2
        + b_1 t
    \right).
\]

The softplus function is

\[
    \operatorname{softplus}(x)
    =
    \log(1 + e^x).
\]

In implementation, softplus is evaluated using the numerically stable form

\[
    \operatorname{softplus}(x)
    =
    \max(x,0) + \log(1 + e^{-|x|}),
\]

which reduces overflow risk for large positive inputs.

The default local-volatility setting uses daily time steps,

\[
    N = 252,\quad
    \sigma_{\min}=0.01,\quad
    a_1=-0.10,\quad
    a_2=0.20,\quad
    b_1=0.00.
\]

The parameter \(a_0\) is chosen so that the at-the-money time-zero volatility
equals \(0.20\):

\[
    a_0
    =
    \log\!\left(
        \exp(0.20-\sigma_{\min}) - 1
    \right).
\]

This local-volatility workload is more demanding than the single-step GBM
case because each path requires repeated state updates and repeated
evaluation of a nonlinear volatility function.  It is therefore used as the
main workload for engine comparison, AD overhead analysis, path-count
scaling, time-step scaling, thread scaling, surface-shape sensitivity, and
local resource-usage analysis.

\subsection{Asian Arithmetic-Average Option}
\label{sec:exp-workload-asian}

The final cloud-profiler experiments also include an Asian call option whose
payoff depends on the arithmetic average of the simulated asset path. In
simplified form, the discounted payoff estimator is

\[
    \hat{P}_{\text{Asian}}
    =
    e^{-rT}
    \frac{1}{M}
    \sum_{i=1}^{M}
    \max\!\left(
        \frac{1}{N}\sum_{j=1}^{N} S_{t_j}^{(i)} - K,\ 0
    \right).
\]

This workload is included because it is path dependent: the payoff depends on
the trajectory rather than only on the terminal value. It therefore introduces
a different memory-access and compilation pattern from the European GBM
baseline. In the final cloud comparison, the Asian workload is evaluated for
supported Python/JAX configurations and contributes to the AD-overhead,
cost-performance, and profiler-selection results.

\section{Engines and Software Stacks}
\label{sec:exp-engines}

Table~\ref{tab:exp-engines} summarises the engine configurations used in the
final experiments.  Engines that are not available in a particular run
are skipped and recorded as unavailable rather than silently excluded.

\begin{table}[htbp]
\centering
\small
\caption{Engine configurations used in the local experimental programme.}
\label{tab:exp-engines}
\begin{tabularx}{\textwidth}{p{0.20\textwidth}p{0.18\textwidth}p{0.18\textwidth}Y}
\toprule
\textbf{Engine label} &
\textbf{Language} &
\textbf{Backend} &
\textbf{Role in evaluation} \\
\midrule

\code{cpu} &
Python &
NumPy / CPU &
Reference vectorised Python implementation.  Used as a baseline software
stack for no-AD CPU throughput and correctness checks. \\

\code{jax} &
Python &
JAX / XLA &
Compiled AD-capable implementation.  Used for no-AD performance,
forward-mode AD, reverse-mode AD, and XLA compiled execution behaviour. \\

\code{cpp} &
C++ &
OpenMP &
Compiled CPU implementation.  Used for low-level CPU throughput and
thread-scaling comparison when the extension module is built successfully. \\

\code{rust} &
Rust &
Rayon &
Compiled CPU implementation.  Used for language/runtime comparison and
parallel CPU throughput on the local-volatility workload. \\
\bottomrule
\end{tabularx}
\end{table}

The exact Python, NumPy, JAX, compiler, operating-system, CPU, memory, and
available backend metadata are captured per run and stored in the database.
This avoids relying on manually written machine descriptions, which can
become incomplete or stale.

\section{Experimental Variables}
\label{sec:exp-variables}

The experiments are designed around a small set of controlled variables.

\subsection{Independent Variables}

The main independent variables are

\begin{itemize}
    \item \textbf{Engine:} NumPy CPU, JAX/XLA, C++/OpenMP, and Rust/Rayon.

    \item \textbf{Path count:} the number of Monte Carlo paths, denoted by
    \(M\).

    \item \textbf{Time steps:} the number of simulated time steps, denoted
    by \(N\).  For the exact GBM workload, \(N=1\) implicitly because the
    terminal distribution is sampled directly.  For local volatility,
    \(N\) is explicitly swept in selected experiments.

    \item \textbf{AD mode:} no AD, forward-mode AD, or reverse-mode AD.
    Forward and reverse modes are evaluated through the JAX engine.

    \item \textbf{Thread count:} requested CPU worker count for strong and
    weak thread-scaling experiments.

    \item \textbf{Volatility-surface preset:} local-volatility surface shape,
    such as flat, equity-skew, or strong-smile, where supported by the
    experiment runner.

    \item \textbf{Cloud instance:} Google Cloud instance type and region for
    final cost-performance and profiler experiments.

    \item \textbf{Profiler policy:} full grid, old probe-based profiler, or
    successive halving.
\end{itemize}

\subsection{Dependent Variables}

The main dependent variables are

\begin{itemize}
    \item mean runtime in milliseconds;
    \item runtime standard deviation, minimum, and maximum;
    \item throughput in paths per second;
    \item peak process memory usage in megabytes;
    \item Monte Carlo price estimate;
    \item Greek estimates for AD-enabled GBM experiments;
    \item absolute and relative error against analytical references where
    available;
    \item AD overhead ratio;
    \item requested and observed process thread counts where available;
    \item coefficient of variation for selected stability experiments.
\end{itemize}

\subsection{Controlled Variables}

Unless explicitly swept, the following are held fixed:

\begin{itemize}
    \item option parameters \(S_0\), \(K\), \(r\), \(\sigma\), and \(T\);
    \item random seed;
    \item workload definition;
    \item local-volatility surface parameters;
    \item warmup count and timed repetition count within each experiment;
    \item database schema and result-writing path;
    \item timing method;
    \item correctness thresholds used for validation.
\end{itemize}

\section{Metrics}
\label{sec:exp-metrics}

For each benchmark cell, the script performs a configurable number of warmup
runs followed by a configurable number of timed runs.  Warmup timings are
discarded.  The timed runs produce

\[
    \bar{t}_{\mathrm{ms}}
    =
    \frac{1}{R}
    \sum_{j=1}^{R} t_j,
\]

where \(R\) is the number of timed repetitions and \(t_j\) is the runtime of
the \(j\)th repetition in milliseconds.

Throughput is reported as

\[
    \mathrm{throughput}
    =
    \frac{M}{\bar{t}_{\mathrm{ms}} \times 10^{-3}}
    \quad
    \text{paths per second}.
\]

For AD experiments, the overhead ratio is computed against a matched no-AD
baseline:

\[
    \mathrm{AD\ overhead}
    =
    \frac{\bar{t}_{\mathrm{AD}}}
         {\bar{t}_{\mathrm{none}}},
\]

where both runs use the same engine, workload, path count, time-step count,
warmup count, and number of timed repetitions.

For the GBM workload, numerical error is measured against the analytical
Black--Scholes reference:

\[
    \mathrm{relative\ error}
    =
    \frac{|x_{\mathrm{MC}} - x_{\mathrm{ref}}|}
         {\max(|x_{\mathrm{ref}}|,\epsilon)},
\]

where \(x\) may be the price, Delta, Vega, or Rho, and \(\epsilon\) is a
small constant used only to avoid division by zero.

For the local-volatility workload, no closed-form analytical reference is
available for the chosen parametric surface.  Cross-engine numerical
consistency is therefore assessed by comparing estimates across engines and
by checking convergence as \(M\) increases.  Where available, the Monte Carlo
standard error is used to interpret whether differences between engines are
consistent with sampling variation.

For selected stability experiments, the coefficient of variation is reported
as

\[
    \mathrm{CV}
    =
    \frac{s_t}{\bar{t}_{\mathrm{ms}}},
\]

where \(s_t\) is the sample standard deviation of the timed repetitions.

\section{Measurement Methodology}
\label{sec:exp-methodology}

\subsection{Warmup and JIT Compilation}

Each script performs warmup repetitions before measurement.  This is
especially important for JAX/XLA, where the first invocation may include
tracing and compilation.  The reported runtime is the steady-state execution
time after warmup, not the cold-start compilation time.

Cold-start costs are not ignored entirely.  Where relevant, they are
discussed separately in the results chapter because they affect short-lived
interactive workloads differently from repeated production-style simulation
runs.

\subsection{Repeated Measurements}

Each benchmark cell is measured over multiple timed repetitions.  The
database stores the mean, standard deviation, minimum, and maximum runtime
rather than a single timing.  This makes it possible to identify unstable
measurements and avoids drawing conclusions from one-off timing noise.

Most main experiments use 7 timed runs and 3 warmup runs.  Smoke tests use
1 timed run and 1 warmup run and are not treated as reported performance
results.  Selected stability experiments use a larger number of timed runs
to estimate measurement variability.

\subsection{Randomness and Correctness}

All workloads use an explicit seed.  Identical seeds improve reproducibility
within each engine, but bit-identical output is not required across engines
because different engines may use different random-number generators.
Cross-engine correctness is therefore assessed using analytical references
where available, and by checking that estimates converge within reasonable
Monte Carlo error as \(M\) increases.

For the European GBM workload, correctness is assessed directly against the
Black--Scholes analytical price and Greeks.  For local volatility, the
emphasis is on cross-engine consistency and convergence with increasing path
count, since the chosen parametric local-volatility model does not have the
same closed-form reference.

\subsection{Fairness Controls}

The following controls are applied across experiments:

\begin{itemize}
    \item Engines receive the same workload parameters and path counts.
    \item AD overhead is computed only against matched no-AD runs.
    \item JAX timings are synchronised so that asynchronous execution does
    not under-report runtime.
    \item Thread-scaling cells are run in isolated subprocesses where
    possible to reduce contamination from persistent thread-pool state.
    \item Failed, skipped, or unavailable engines are recorded explicitly
    rather than removed silently from the analysis.
    \item Smoke-test results are used only to validate execution and are not
    interpreted as performance results.
\end{itemize}

\subsection{Data Exclusion Policy}

Runs are excluded from the final comparative analysis only under predefined
conditions:

\begin{itemize}
    \item the run terminates with an exception;
    \item the detected engine or backend does not match the intended
    experiment cell;
    \item the result violates the experiment's correctness threshold;
    \item required metadata such as engine, workload, path count, time-step
    count, or runtime is missing;
    \item the machine is known to have been under unrelated heavy load during
    the run.
\end{itemize}

Excluded runs are not deleted from the database.  They remain available for
audit, but are filtered out of the final completed-run queries.

\section{Experiment Scripts}
\label{sec:exp-scripts}

Table~\ref{tab:exp-scripts} summarises the core experiment scripts used in
the final evaluation.

\begin{table}[htbp]
\centering
\small
\caption{Core experiment scripts and their roles.}
\label{tab:exp-scripts}
\begin{tabularx}{\textwidth}{p{0.28\textwidth}p{0.23\textwidth}Y}
\toprule
\textbf{Script} &
\textbf{Experiment type or labels} &
\textbf{Purpose} \\
\midrule

\code{run_european_baseline.py} &
\code{european_baseline} &
Compares no-AD throughput across available engines on the single-step
European GBM workload. \\

\code{run_european_ad_analysis.py} &
\code{european_ad_analysis} &
Measures JAX forward- and reverse-mode AD overhead and validates Greek
accuracy against Black--Scholes references. \\

\code{run_european_scalability.py} &
\code{european_scalability} &
Sweeps path count \(M\) to measure how runtime and throughput scale on the
simple GBM workload. \\

\code{run_thread_scalability.py} &
\code{thread_scalability_strong};
\code{thread_scalability_weak} &
Measures strong and weak CPU thread scaling on the GBM workload. \\

\code{run_localvol_suite.py} &
\code{LV0}--\code{LV10} &
Runs the local-volatility benchmark suite, including engine comparison,
AD-overhead scaling, path-count scaling, time-step scaling, thread scaling,
surface-shape sensitivity, timing stability, and local stress testing. \\

\code{run_profiler_vs_grid.py} &
\code{sha\_cloud\_profiler\_v4} &
Runs the final cloud full-grid, old-profiler, and SHA comparison, recording
cloud metadata, cost estimates, Pareto-frontier recovery, regret, and
promotion decisions. \\
\bottomrule
\end{tabularx}
\end{table}

\section{Individual Experiment Designs}
\label{sec:exp-designs}

\subsection{European GBM Baseline}
\label{sec:exp-baseline}

The European GBM baseline establishes a simple no-AD performance reference
and a correctness check against the Black--Scholes analytical price.

\paragraph{Script.}
\code{experiments/run_european_baseline.py}

\paragraph{Workload.}
European GBM with fixed option parameters.

\paragraph{Engines.}
NumPy CPU, JAX/XLA, C++/OpenMP, and Rust/Rayon where available.

\paragraph{Purpose.}
This experiment provides a low-cost first comparison across engines and
checks that the framework can run each implementation through the same
configuration and database-writing path.

\paragraph{Primary outputs.}
Mean runtime, runtime standard deviation, throughput, price estimate,
relative price error, memory usage, and environment metadata.

\subsection{European GBM AD Overhead and Greek Accuracy}
\label{sec:exp-ad}

This experiment isolates the marginal cost of automatic differentiation on
the analytical GBM workload.

\paragraph{Script.}
\code{experiments/run_european_ad_analysis.py}

\paragraph{Workload.}
European GBM.

\paragraph{Engine.}
JAX/XLA.

\paragraph{Sweep.}
\[
    M \in
    \{1{,}000,\ 5{,}000,\ 10{,}000,\ 50{,}000,\ 100{,}000\},
    \qquad
    \text{AD mode}
    \in
    \{\text{none},\ \text{forward},\ \text{reverse}\}.
\]

\paragraph{Purpose.}
The experiment compares no-AD execution with forward- and reverse-mode AD
under identical workload settings.  It also checks whether AD-computed
Greeks are numerically consistent with the analytical Black--Scholes Delta,
Vega, and Rho.

\paragraph{Primary outputs.}
AD overhead ratio, Greek estimates, Greek errors, runtime, throughput, and
memory usage.

\subsection{European GBM Path-Count Scalability}
\label{sec:exp-path-scaling}

This experiment measures how each engine behaves as the number of Monte
Carlo paths increases on the simple GBM workload.

\paragraph{Script.}
\code{experiments/run_european_scalability.py}

\paragraph{Workload.}
European GBM.

\paragraph{Sweep.}
\[
    M \in
    \{1{,}000,\ 5{,}000,\ 10{,}000,\ 50{,}000,\ 100{,}000\}.
\]

\paragraph{Purpose.}
The experiment identifies fixed overheads at small path counts and
throughput behaviour at larger path counts.  It is used as a simple
reference before interpreting the more expensive local-volatility workload.

\paragraph{Primary outputs.}
Runtime, throughput, price error, and memory usage as functions of \(M\).

\subsection{European GBM Thread Scalability}
\label{sec:exp-thread-scaling}

The GBM thread-scaling experiment measures how engines use local CPU
parallelism on a simple workload.

\paragraph{Script.}
\code{experiments/run_thread_scalability.py}

\paragraph{Workload.}
European GBM.

\paragraph{Engines.}
JAX/XLA and C++/OpenMP, with Rust/Rayon included where supported.

\paragraph{Regimes.}
Two scaling regimes are evaluated:

\begin{itemize}
    \item \textbf{Strong scaling:} \(M\) is fixed while requested thread
    count varies.  Ideal strong scaling means runtime decreases in
    proportion to the number of threads.

    \item \textbf{Weak scaling:} \(M\) increases in proportion to requested
    thread count.  Ideal weak scaling means runtime remains approximately
    constant as total work and total resources increase together.
\end{itemize}

\paragraph{Primary outputs.}
Runtime, throughput, requested thread count, observed thread counts, and
scaling efficiency.

\subsection{Local-Volatility Benchmark Suite}
\label{sec:exp-localvol-suite}

The local-volatility benchmark suite is the main experimental programme in
the final evaluation.

\paragraph{Script.}
\code{experiments/run_localvol_suite.py}

\paragraph{Workload.}
European local volatility with default daily discretisation \(N=252\),
unless \(N\) is explicitly swept.

\paragraph{Engines.}
NumPy CPU, JAX/XLA, C++/OpenMP, and Rust/Rayon where available.  AD modes are
evaluated through JAX.

\paragraph{Experiment groups.}
The suite contains the following experiment groups:

\begin{itemize}
    \item \textbf{LV-0 Smoke test:} checks that all available engines can run
    the local-volatility workload before expensive measurements are made.
    This experiment is not reported as a performance result.

    \item \textbf{LV-1 Engine and \(M\)-scalability:} compares no-AD
    throughput across NumPy, JAX, C++, and Rust for
    \(M \in
    \{10{,}000,25{,}000,50{,}000,100{,}000,250{,}000,500{,}000\}\)
    at \(N=252\).

    \item \textbf{LV-2 JAX AD \(M\)-scaling:} compares no AD,
    forward-mode AD, and reverse-mode AD in JAX for
    \(M \in
    \{5{,}000,10{,}000,25{,}000,50{,}000,100{,}000,250{,}000\}\)
    at \(N=252\).

    \item \textbf{LV-3 Time-step scaling:} sweeps the number of time steps
    \(N\) at fixed \(M\) to separate path-count scaling from
    discretisation-cost scaling.

    \item \textbf{LV-4 JAX AD time-step scaling:} measures how AD overhead
    changes as the number of time steps increases.

    \item \textbf{LV-5 Memory pressure:} records peak process memory usage
    for larger local-volatility runs, with particular attention to JAX AD
    modes.

    \item \textbf{LV-6 Strong thread scaling:} evaluates whether fixed-size
    local-volatility workloads speed up when the requested CPU thread count
    is increased over \(1,2,4\).

    \item \textbf{LV-7 Weak thread scaling:} evaluates whether runtime
    remains stable when both requested thread count and total path count
    increase together.

    \item \textbf{LV-8 Surface-shape sensitivity:} compares alternative
    local-volatility surface presets, such as flat, equity-skew, and
    strong-smile, where supported by the runner.

    \item \textbf{LV-9 Timing stability:} repeats selected key cells with a
    larger number of timed runs to estimate measurement variability.

    \item \textbf{LV-10 Stress limit:} tests the largest practical
    local-volatility configurations that fit within the 4-core, 16 GB RAM
    Codespace environment.
\end{itemize}

\paragraph{Purpose.}
This suite tests whether the performance conclusions from the simple GBM
workload still hold for a more realistic MC+AD workload.  It provides the
primary evidence for programming-language/runtime comparison, AD overhead,
path-count scalability, time-step scalability, thread scalability, local
resource usage, and measurement stability.

\paragraph{Primary outputs.}
Mean runtime, runtime standard deviation, throughput, price estimate,
AD overhead ratio, memory usage, requested and observed thread counts, and
coefficient of variation for selected stability runs.

\section{Final Cloud and Profiler Experiments}
\label{sec:exp-cloud-profiler}

The final experimental extension evaluates cloud cost-performance and the
budget-aware profiler. It does not replace the earlier local experiments:
those experiments provide the detailed workload, AD, scaling, and resource
analysis. The cloud experiment instead asks whether the completed framework can
use stored metadata and cheap probes to make cost-aware full-benchmark
selection decisions.

\subsection{Full-Grid Cloud Baseline}
\label{sec:exp-full-grid-cloud}

The headline full-grid baseline is the \texttt{sha\_cloud\_profiler\_v4}
experiment. It is recorded at git commit \texttt{13789b8} on a Google Cloud
\texttt{n2-standard-4} instance in \texttt{europe-west1-b}, with region
\texttt{europe-west1}. The full-grid baseline contains 16 supported candidate
configurations and three full path counts,
\[
    M \in \{10{,}000,\ 50{,}000,\ 100{,}000\},
\]
giving 48 full benchmark cells.

The candidate configurations span the European, Asian, and local-volatility
workloads, supported no-AD engines, and JAX AD modes. Unsupported combinations,
such as AD for compiled no-AD-only engines, are excluded before the grid is
formed. The full grid provides the reference decisions against which profiler
quality is judged.

\subsection{Old Profiler Baseline}
\label{sec:exp-old-profiler}

The old profiler performs cheap probe runs, ranks candidate configurations, and
selects a conservative subset for full benchmarking. In the final experiment it
selects 30 of the 48 full-grid cells. Its purpose is to provide a stronger
baseline than ``no profiling'': SHA must be compared not only with exhaustive
search, but also with the earlier probe-selection strategy.

\subsection{Successive Halving Profiler}
\label{sec:exp-sha-profiler}

The SHA profiler evaluates all candidates at low path counts and promotes a
smaller active set to larger probe budgets. The recorded progression uses probe
path counts including \(M=1{,}000\), \(M=5{,}000\), and \(M=25{,}000\). After
the promotion rounds and conservative safeguards, SHA selects 24 of the 48
full-grid cells for full benchmarking.

The SHA result is evaluated using the same decision-quality criteria as the old
profiler: full runs saved, Pareto recovery, runtime regret, cost regret,
Spearman probe-to-full rank correlation, and probe path-budget saving. This
ensures that the profiler is not credited merely for reducing work; it must
also preserve the decisions supported by the full grid.

\section{Execution Order}
\label{sec:exp-execution-order}

The experiments are run in an order that first validates correctness on a
simple analytical workload and then moves to the more expensive
local-volatility suite and final cloud-profiler evaluation.

\begin{enumerate}
    \item \textbf{European GBM validation:} run the baseline GBM workload and
    confirm that completed engines produce prices close to the
    Black--Scholes analytical reference.

    \item \textbf{European AD validation:} run the JAX forward- and
    reverse-mode AD experiment and compare Delta, Vega, and Rho against
    Black--Scholes references.

    \item \textbf{European path-count and thread scaling:} run the simpler
    GBM scalability experiments to establish low-cost baseline behaviour.

    \item \textbf{Local-volatility smoke test:} run LV-0 to confirm that all
    available engines can execute the multi-step workload.

    \item \textbf{Local-volatility engine and AD sweeps:} run LV-1 and LV-2
    to collect the main local-volatility performance and AD-overhead data.

    \item \textbf{Local-volatility \(N\)-scaling:} run LV-3 and LV-4 to
    measure how runtime and AD overhead change with the number of time
    steps.

    \item \textbf{Local-volatility resource and parallelism experiments:}
    run LV-5 to LV-10 to evaluate memory usage, thread scaling,
    surface-shape sensitivity, timing stability, and practical local stress
    limits.

    \item \textbf{Cloud full-grid and profiler comparison:} run
    \texttt{sha\_cloud\_profiler\_v4} to compare full-grid benchmarking with
    the old profiler and SHA under recorded cloud cost metadata.
\end{enumerate}

This order keeps the analytically validated GBM workload as a correctness
anchor, while making the local-volatility workload the main performance
benchmark for detailed scaling analysis and the cloud-profiler experiment the
main evidence for cost-aware selection.

\section{Database Queries for Analysis}
\label{sec:exp-analysis-queries}

The final analysis uses database queries over completed runs.  The exact
query predicates may be adjusted to match the experiment labels stored by
the runner, but the following views define the main result tables.

\subsection{European Baseline Throughput}

\begin{lstlisting}[style=sqlstyle,language=SQL]
SELECT engine,
       mean_runtime_ms,
       std_runtime_ms,
       throughput_paths_per_sec,
       result_value,
       rel_price_error,
       memory_peak_mb
FROM runs
WHERE experiment_type = 'european_baseline'
  AND status = 'completed'
ORDER BY throughput_paths_per_sec DESC;
\end{lstlisting}

\subsection{European AD Overhead and Greek Accuracy}

\begin{lstlisting}[style=sqlstyle,language=SQL]
SELECT M,
       ad_mode,
       mean_runtime_ms,
       ad_overhead_ratio,
       greek_delta,
       analytical_delta,
       rel_delta_error,
       greek_vega,
       analytical_vega,
       rel_vega_error,
       greek_rho,
       analytical_rho,
       rel_rho_error
FROM runs
WHERE experiment_type = 'european_ad_analysis'
  AND status = 'completed'
ORDER BY M, ad_mode;
\end{lstlisting}

\subsection{Local-Volatility Engine Scaling}

\begin{lstlisting}[style=sqlstyle,language=SQL]
SELECT experiment_type,
       engine,
       M,
       N,
       mean_runtime_ms,
       std_runtime_ms,
       throughput_paths_per_sec,
       result_value,
       memory_peak_mb
FROM runs
WHERE experiment_type LIKE 'localvol%'
  AND ad_mode = 'none'
  AND status = 'completed'
ORDER BY experiment_type, M, engine;
\end{lstlisting}

\subsection{Local-Volatility AD Overhead}

\begin{lstlisting}[style=sqlstyle,language=SQL]
SELECT experiment_type,
       M,
       N,
       ad_mode,
       mean_runtime_ms,
       std_runtime_ms,
       ad_overhead_ratio,
       throughput_paths_per_sec,
       memory_peak_mb,
       result_value
FROM runs
WHERE experiment_type LIKE 'localvol%'
  AND engine = 'jax'
  AND status = 'completed'
ORDER BY experiment_type, N, M, ad_mode;
\end{lstlisting}

\subsection{Thread Scaling}

\begin{lstlisting}[style=sqlstyle,language=SQL]
SELECT engine,
       experiment_type,
       requested_threads,
       observed_threads_max,
       M,
       N,
       mean_runtime_ms,
       throughput_paths_per_sec
FROM runs
WHERE experiment_type LIKE '%thread%'
  AND status = 'completed'
ORDER BY experiment_type, engine, M, N, requested_threads;
\end{lstlisting}

\section{Reproducibility}
\label{sec:exp-reproducibility}

Each completed run stores the full workload configuration and environment
metadata needed to reproduce or audit the measurement.  The database record
includes

\begin{itemize}
    \item the serialised workload configuration;
    \item a configuration hash where available;
    \item engine, language, backend, and AD mode;
    \item path count, time-step count, option parameters, and seed;
    \item local-volatility surface parameters or preset;
    \item timing statistics;
    \item price and Greek outputs where available;
    \item analytical reference values for GBM runs;
    \item numerical error metrics where a reference is available;
    \item memory usage;
    \item requested and observed thread counts where available;
    \item CPU, operating-system, Python, NumPy, and JAX metadata.
\end{itemize}

This metadata ensures that every reported number can be traced back to a
specific workload, software stack, and hardware environment.  The experiment
database is append-only during collection: failed, skipped, and completed
runs are all preserved so that the final analysis does not depend on
unrecorded manual filtering.

\section{Scope of Excluded Extensions}
\label{sec:exp-extensions}

The original project proposal included cloud resource selection, GPU
acceleration, distributed execution, mixed precision, BLAS backend sweeps,
and vendor-claim validation through emerging languages such as Mojo. Cloud
resource selection is partially addressed through the final cost-performance
and profiler experiments. GPU acceleration, distributed execution, mixed
precision, broader BLAS-environment sweeps, and Mojo vendor-claim validation
remain future extensions.

The final results therefore combine local CPU benchmarking across NumPy,
JAX/XLA, C++/OpenMP, and Rust/Rayon with a targeted cloud-profiler study.
The local campaign provides detailed analysis of local-volatility
performance, AD overhead, path-count scaling, time-step scaling, thread
scaling, and measurement stability. The cloud campaign provides measured
cost-performance, Pareto-frontier, and profiler-selection evidence for the
completed framework.

\section{Summary}
\label{sec:exp-summary}

The experimental setup defines a controlled methodology for MC+AD workloads,
combining a 4-core, 16 GB RAM Codespace campaign with final Google Cloud
profiler experiments.  The European GBM workload provides analytical
validation and a simple baseline. The parametric local-volatility workload is
the main detailed performance benchmark, adding multi-step simulation,
nonlinear model structure, and higher AD complexity. The Asian workload adds a
path-dependent case to the cloud-profiler grid.

Together, the experiments evaluate runtime performance, scalability,
automatic differentiation overhead, numerical consistency, thread behaviour,
local resource usage, cloud cost-performance, and budget-aware profiler
quality across NumPy, JAX/XLA, C++/OpenMP, and Rust/Rayon where supported.
The next chapter reports the measurements obtained from these experiments
and discusses how far they answer the project research questions.
